\documentclass[14px]{article}
\usepackage{xeCJK}
\usepackage[frenchb]{babel}
\usepackage[T1]{fontenc}
\usepackage[utf8]{inputenc}
\usepackage{textcomp}
\usepackage{amssymb}
\usepackage[ruled,longend]{algorithm2e}
\usepackage{amsmath}
\usepackage{latexsym}
\usepackage{fancyhdr}
\usepackage{geometry}
\usepackage{setspace}
\usepackage[colorlinks,linkcolor=blue]{hyperref}
% Image
\usepackage{graphicx}
\usepackage{float}
\usepackage{subfigure}
\usepackage{enumerate}

\begin{document}
\setlength{\parindent}{0pt}
\begin{titlepage}
	\begin{center}
		% Upper part of the page
		\includegraphics[width=0.35\textwidth]{logo.png}\\[1cm]
		\textsc{\Large Rapport de PC3R}\\[0.5cm]
		% Title
		{ \huge \bfseries Système de Gestion de Location}\\[0.4cm]
		% Author and supervisor		
		\begin{minipage}{0.4\textwidth}
			\begin{flushleft} \large
				\emph{Author:}\\
				Qiwei \textsc{XIAN}\\
				Ruiwen \textsc{WANG}\\
			\end{flushleft}
		\end{minipage}
		\begin{minipage}{0.4\textwidth}	
				\begin{flushright} \large
				\emph{Professeur:} \\
				Prof.\textsc{Romain Demangeon}
			\end{flushright}
		\end{minipage}
		\vfill
		% Bottom of the page
		{\large \today}
	\end{center}
\end{titlepage}
\clearpage

\tableofcontents
\thispagestyle{empty}
\clearpage
\section{Introduction}
Système de gestion de location est une plateforme communautaire de location et de réservation de logements personnels. Les fonctionnalités du système est ressemble à Airbnb. Il permet aux utilisateurs de louer leur propriétés immobilières et de réserver un logement de l'autre utilisateur.

Chaque l'utilisateur doit inscrire sur le système afin d'obtenir un compte, il bénéficie de la service du système en utilisant ce compte. Il peut créer une espace de client et enregistrer les informations publiques, chaque utilisateur peut consulter les informations 
des autres utilisateurs.

Il peut tourver des logements qui satisfait à ses beosins, par exemple sous la condition de période de location, la ville, le nombre des locataires, ainsi que le droit de fumer.
L'utilisateur peut aussi gérer leur propriétés par ce systèmes, il d'abord ajoute ses logements à louer dans le compte et met les contraintes pour les locataires, par exemple le nombre des locataires ou l'interdiction de fumer, etc.  En plus il peut les modifier ou supprimer comme il veut. Lors que l'utilisateur reçoit les demandes envoyées par les autres utilisateurs. Il a droit de la refuser ou accepter, mais si la demande a risque de causer un conflit, le système va le faire remarquer au propriétaire.
 
\section{API Web choisie}

\section{Fonctionnalités}
Le système est séparer en trois modules, \textbf{Account}, \textbf{Demand}, \textbf{Resource}. \\
Le module \textbf{Account} sert à gérer les données de l'authentification et l'espace de l'utilisateur. Le module \textbf{Resource} permet de gérer les données de propriété. Et le module \textbf{Demand} sert à gérer les données des demands.\\
Cela permet de faciliter le développement parallèle et la maintenance de l'application. On a réaliser les fonctionnalités de l'ajoute, de la suppression et de la consultation, ainsi que la modification de ces trois modules. 

\section{Cas d'utilisations}
\subsection{Les fonctionnalités du module Compte}
\subsubsection{Inscription}
Acteur : L'utilisateur\\
Contexte : L'utilisateur crée un compte.\\
Scénario principal: 
\begin{enumerate}
	\item L'utilisateur entre dans le page d'inscription.
	\item L'utilisateur saisit l'identifiant et le mot de passe.
	\item L'utilisateur clique Signup pour soumettre les information au serveur.
	\item Le système répond le message de success et afficher sur l'écran.	
\end{enumerate}
Cas particulier:\\
4a. L'identifiant est déjà utilisé dans le système, le système affiche le message d'erreur.

\subsubsection{Authentification}
Acteur: L'utilisateur\\
Contexte: L'utilisateur s'authentifie dans le système par son compte.
Scénario principal:
\begin{enumerate}
	\item L'utilisateur entre dans le page de l'authentification.
	\item L'utilisateur saisit l'identifiant et le mot de passe.
	\item L'utilisateur clique le button Login pour soumettre les informations au serveur.
	\item Le serveur vérifie l'identifiant et le mot de passe.
	\item Le serveur rend le page web de \textbf{mainPage.jsp}, l'identifiant est stocké dans la session.
\end{enumerate}
Cas particulier: \\
4a. L'identifiant n'est pas reconnu, le système affiche le message d'erreur.\\
4b. Le mot de passe n'est pas correspondant à l'identifiant, le système affiche le message d'erreur.

\subsubsection{Consulter l'espace du client}
Acteur: L'utilisateur\\
Contexte: L'utilisateur connecté veut consulter son espace du client.
Scénario principal:
\begin{enumerate}
	\item L'utilisateur clique le button Profile dans mainPage.
	\item Le serveur génère dynamiquement l'espace personnelle de l'utilisateur.   
\end{enumerate}
cas particulier:\\
2a. L'utilisateur n'a pas encore son espace du client, le serveur génère le page de creation d'espace, permet à l'utilisateur de créer son espace.


\subsubsection{Créer l'espace du client}
Acteur: L'utilisateur\\
Contexte: L'utilisateur consulte sont espace première fois
Scénario principal:
\begin{enumerate}
	\item L'utilisateur clique le button Profile première fois dans le mainPage.
	\item Le serveur repond un message et génère dynamiquement le formulaire pour créer l'espace du client.
	\item L'utilisateur saisit les informations publiques.
	\item L'utilisateur clique le button Créer afin de soumettre les information.
	\item Le serveur crée l'espace du client et le stocke dans la base de données, ainsi que génère le page web de l'espace du client.
\end{enumerate}

\subsubsection{Modifier l'espace du client}
Acteur: L'utilisateur
Contexte: L'utilisateur connecté se trouve dans son espace du client.
Scénario principal:
\begin{enumerate}
	\item L'utilisateur clique le button \textbf{here} afin de modifier son 
	\item Le serveur génère le page permet au client de modifier son espace.
\end{enumerate}

\subsection{Les fonctionnalités du module Ressource}
\subsubsection{Consulter les ressources}
Acteur: L'utilisateur\\
Contexte: L'utilisateur s'authentifié se trouve à mainPage.
Scénario principal:
\begin{enumerate}
	\item L'utilisateur clique le button \textbf{Your house}.
	\item Le serveur cherche dans la base de donnée et génère dynamiquement le page sert à afficher les informations de propriétés.
\end{enumerate}

\subsubsection{Ajouter la ressource}
Acteur: L'utilisateur\\
Contexte: L'utilisateur connecté se trouve dans le page \textbf{Your House}.
Scénario principal:
\begin{enumerate}
	\item L'utilisateur clique le button \textbf{Add House}.
	\item La page est générée dynamiquement par js permet à l'utilisateur de saisir les informations de la ressource.
	\item L'utilisateur clique le button \textbf{Créer} pour soumettre.
	\item Le serveur vérifie que le status de l'authentification de l'utilisateur est connecté.
	\item Le serveur stocke les informations de la ressource.
	\item Le serveur affiche un message de success sur l'écran.
\end{enumerate}



\section{Stockage des données}

\section{Structure du serveur}

\section{Structure du client}

\section{Requêtes et réponse}

\section{Conclusion}

	
\end{document}